\documentclass[parskip=full]{scrreprt}

\RedeclareSectionCommand[pagestyle=plain,afterskip=10pt plus 1pt]{chapter}
\setkomafont{chapter}{\mdseries\headingfont\fontsize{40}{40}\selectfont\color{black!80}}
\setkomafont{pageheadfoot}{\normalsize}

\def\pnumbox#1{#1\hspace*{7cm}}
\def\gobble#1{}
\DeclareTOCStyleEntry[
  indent=0pt,
  beforeskip=0pt,
  entryformat=\itshape,
  entrynumberformat=\textcolor{oldred},
  numwidth=2em,
  linefill=\hfill,
  pagenumberbox=\pnumbox,
  pagenumberformat=\itshape
]{tocline}{part}

% use commented values if there are no parts
\DeclareTOCStyleEntry[
  indent=0pt,
  beforeskip=-.5\baselineskip,
  entrynumberformat=\textcolor{oldred},
  numwidth=3em,
  linefill=\hfill,
  pagenumberbox=\pnumbox
]{tocline}{section}

\DeclareTOCStyleEntry[
  indent=0pt,
  beforeskip=-\parskip,
  entrynumberformat=\gobble,
  entryformat=\ltseries,
  numwidth=2em,
  linefill=\hfill,
  pagenumberbox=\pnumbox,
  pagenumberformat=\ltseries
]{tocline}{subsection}

\usepackage{polyglossia}
\setdefaultlanguage{english}

\usepackage{fontspec}

\setmainfont{Source Sans Pro}[
  ItalicFont = Source Sans Pro Italic,
  BoldFont = Source Sans Pro Bold,
  BoldItalicFont = Source Sans Pro Bold Italic,
  FontFace = {lt}{n}{Source Sans Pro Light},
  FontFace = {lt}{it}{Source Sans Pro Light Italic},
  FontFace = {sb}{n}{Source Sans Pro Semibold},
  FontFace = {sb}{it}{Source Sans Pro Semibold Italic},
  Numbers = Proportional,
  Ligatures = Common
]
\DeclareRobustCommand{\ltseries}{\fontseries{lt}\selectfont}
\DeclareRobustCommand{\sbseries}{\fontseries{sb}\selectfont}
\DeclareTextFontCommand{\textlt}{\ltseries}
\DeclareTextFontCommand{\textsb}{\sbseries}

\newfontfamily\headingfont{Fredericka the Great}

\usepackage[pass]{geometry}
\newgeometry{twoside,inner=20mm,outer=40mm,top=20mm,bottom=40mm}

\usepackage{url}
\urlstyle{same}

\usepackage{microtype}
\microtypesetup{verbose=silent}

\usepackage{booktabs,array,longtable}
\newcolumntype{L}[1]{>{\raggedright\let\newline\\\arraybackslash\hspace{0pt}}p{#1}}

\usepackage{graphicx}

\usepackage{xcolor}
\definecolor{oldred}{rgb}{.8313,0,0}

\usepackage{pdfpages}

\usepackage{scrlayer-scrpage}
\pagestyle{scrheadings}
\clearscrheadfoot
\cfoot[\thepage]{\thepage}
\pagenumbering{roman}

\usepackage{enumitem}
\setlist[description]{%
  style=nextline,%
  leftmargin=2em,%
  first=\ltseries,%
  font=\normalfont%
}
\def\lyrefitem#1{\item[\lyref{#1}]}

\usepackage{multicol}

\makeatletter

\def\@firstofthree#1#2#3{#1}
\def\@secondofthree#1#2#3{#2}
\def\@thirdofthree#1#2#3{#3}
\def\Dotfill{\leavevmode\cleaders\hb@xt@ .75em{\hss .\hss }\hfill \kern \z@}

\def\lyrefnumber#1{\expandafter\@setref\csname r@#1\endcsname\@firstofthree{#1}}
\def\lyreftitle#1{\expandafter\@setref\csname r@#1\endcsname\@secondofthree{#1}}
\def\lyrefpage#1{\expandafter\@setref\csname r@#1\endcsname\@thirdofthree{#1}}

\def\lyref#1{%
  \begingroup%
  \makebox[0pt][l]{\color{oldred}\lyrefnumber{#1}}\hspace*{2em}%
  \lyreftitle{#1}\Dotfill\lyrefpage{#1}%
  \endgroup%
}
\InputIfFileExists{../out/lilypond.ref}{}{\InputIfFileExists{../lilypond.ref}{}{}}


\newcommand\fancytitlehead{
	\headingfont%
	\fontsize{80}{80}\selectfont\textcolor{black!80}{\@ifundefined{@shortname}{\@lastname}{\@shortname}.}\\[15pt]%
	\fontsize{60}{60}\selectfont\@ifundefined{@shorttitle}{\@title}{\@shorttitle}.%
}


\def\firstname#1{\def\@firstname{#1}}
\def\lastname#1{\def\@lastname{#1}}
\def\shortname#1{\def\@shortname{#1}}
\def\shorttitle#1{\def\@shorttitle{#1}}
\def\namesuffix#1{\def\@namesuffix{#1}}
\def\instrumentation#1{\def\@instrumentation{#1}}
\def\parts#1{\def\@parts{#1}}

\firstname{\relax}
\lastname{\relax}
\namesuffix{\relax}
\instrumentation{\relax}
\parts{\relax}


\def\maketitle{%
\begin{titlepage}%
	\Large%
	{\@titlehead}%
	\vfill%
	{\strut\@firstname}\\%
	{\sbseries\color{oldred}\strut\@lastname}\\%
	{\strut\@namesuffix}%
	\vfill%
	{\sbseries\@title}\\%
	{\@subtitle}\\[\baselineskip]%
	{\itshape\@instrumentation}%
	\vfill%
	{\itshape\@parts}\hspace*{\fill}\raisebox{0pt}[0pt][0pt]{\includegraphics{ees_logo}}%
\end{titlepage}%
}


\newif\iftemplate\templatetrue
\newif\ifprintreport\printreportfalse
\directlua{
scores = {
  vl1 = "Violino I",
  vl2 = "Violino II",
  vla = "Viola",
  coro = "Coro",
  org = "Organo",
  b = "Bassi",
  full_score = "Full Score"
}

last_arg = arg[\string#arg]
texio.write("Last argument: " .. last_arg)
if not (scores[last_arg] == nil) then
  tex.print("\string\\def\string\\lypdfname{" .. last_arg .. "}")
  tex.print("\string\\parts{" .. scores[last_arg] .. "}")
  if (last_arg == "full_score") then
    tex.print("\string\\printreporttrue")
  end
end
}

% uncomment if the full score is in landscape format
\ifprintreport
\geometry{landscape,outer=127mm}
\fi

\@ifundefined{lypdfname}{%
  \templatefalse
  \printreporttrue
  \parts{Draft}
}{\templatetrue}

\makeatother






\begin{document}
\frenchspacing

\titlehead{\fancytitlehead}
\firstname{Anton Adam}
\lastname{Bachschmid}
\title{Stabat mater}
\subtitle{(D-Eu Esl II 60)}
\instrumentation{S, A, T, B (solo), S, A, T, B (coro), 2 vl, vla, b, org}
\maketitle


\thispagestyle{empty}

\vspace*{\fill}

\raisebox{-4mm}{\includegraphics{byncsaeu}}\hspace*{1em}Wolfgang Esser-Skala, 2020

© 2020 by Wolfgang Esser-Skala. This edition is licensed under the Creative Commons Attribution-NonCommercial-ShareAlike 4.0 International License. To view a copy of this license, visit \url{http://creativecommons.org/licenses/by-nc-sa/4.0/}. 

Music engraving by LilyPond 2.18.0 (\url{http://www.lilypond.org}).\\
Front matter typeset with Source Sans Pro and Fredericka the Great.

\textit{First version, July 2020}

\vspace*{2cm}

\ifprintreport
\chapter*{Critical Report.}

This edition bases upon a copy in the Universitätsbibliothek Eichstätt-Ingolstadt. The digital version of the manuscript is available at \url{https://nbn-resolving.org/urn:nbn:de:bvb:824-esl-ii-60-4} (siglum Esl II 60).

In general, this edition closely follows the manuscript. Any changes that were introduced by the editor are indicated by italic type (lyrics, dynamics and directions), parentheses (expressive marks and bass figures) or dashes (slurs and ties). Accidentals are used according to modern conventions. Asterisks denote changes that are clarified in the detailed remarks below.\footnote{Abbreviations: A, alto; B, bass; b, basses; Ms, manuscript; org, organ; r, rest; S, soprano; T, tenor; vl, violin; vla, viola.}

\bigskip

\begin{longtable}{lll L{10cm}}
	\toprule
	\itshape Mov. & \itshape Bar & \itshape Staff & \itshape Note \\
	\midrule \endhead
	V–VII & 44      & vl 2 & 5th eighth in Ms: d′8 \\
	VIII  & 40      & org  & 2nd quarter in Ms: e16–g16–e16–g16 \\
	XX    & 14      & vla  & 13th sixteenth in Ms: r16 \\
	      & 39–40   & org  & bars in Ms: c″+es″1 and fis′+c″1 \\
	      & 48–50   & org  & upper voice in Ms: c″1–d″1–g′1 \\
	      & 81–83   & org  & bars in Ms: es2–g′+b′2, as′+c″1, and d′+f″1 \\
	      & 92–95   & org  & in Ms unison with S and A \\
	      & 109–110 & org  & upper voice missing in Ms \\
	\bottomrule
\end{longtable}


This edition has been compiled and checked with utmost diligence. Nevertheless, errors and mistakes cannot be totally excluded. Please report any errors and mistakes to \url{wolfgang@esser-skala.at} or create an issue or pull request on the edition’s GitHub page \url{https://github.com/skafdasschaf/bachschmid-stabat-mater}. Your help will be greatly appreciated.

\bigskip
\textit{Salzburg, July 2020\\
Wolfgang Esser-Skala}

\cleardoublepage
\chapter*{Lyrics.}

\bigskip
\begin{multicols}{2}
\begin{enumerate}[label=\textcolor{oldred}{\Roman*},leftmargin=0pt]
\item Stabat mater dolorosa\\
iuxta crucem lacrimosa,\\
dum pendebat filius.

\item Cuius animam gementem,\\
contristantem et dolentem\\
pertransivit gladius.

\item O quam tristis et afflicta\\
fuit illa benedicta\\
mater unigeniti!

\item Quae maerebat et dolebat,\\
pia mater cum videbat\\
nati poenas incliti.

\item Quis est homo, qui non fleret,\\
matrem Christi si videret\\
in tanto supplicio?

\item Quis non posset contristari,\\
piam matrem contemplari\\
dolentem cum filio?

\item Pro peccatis suae gentis\\
vidit Iesum in tormentis\\
et flagellis subditum.

\item Vidit suum dulcem natum\\
morientem, desolatum,\\
dum emisit spiritum.

\item Eia, mater, fons amoris,\\
me sentire vim doloris\\
fac, ut tecum lugeam.

\item Fac, ut ardeat cor meum\\
in amando Christum Deum,\\
ut sibi complaceam.

\item Sancta mater, istud agas,\\
crucifixi fige plagas\\
cordi meo valide.

\item Tui nati vulnerati,\\
tam dignati pro me pati,\\
poenas mecum divide.

\item Fac me tecum pie flere,\\
crucifixo condolere,\\
donec ego vixero

\item Iuxta crucem tecum stare,\\
te libenter sociare\\
in planctu desidero.

\item Virgo virginum praeclara,\\
mihi iam non sis amara,\\
fac me tecum plangere.

\item Fac, ut portem Christi mortem,\\
passionis fac consortem\\
et plagas recolere.

\item Fac me plagis vulnerari,\\
cruce hac inebriari\\
ob amorem filii.

\item Inflammatis et accensus\\
per te, virgo, sim defensus\\
in die iudicii.

\item Fac me cruce custodiri,\\
morte Christi praemuniri,\\
confoveri gratia.

\item Quando corpus morietur,\\
fac ut anima donetur\\
paradisi gloriae.
\end{enumerate}
\end{multicols}

\cleardoublepage
\chapter*{Contents.}

\vspace*{\baselineskip}
\InputIfFileExists{../out/lilypond.toc}{}{\InputIfFileExists{../lilypond.toc}{}{}}

\cleardoublepage
\fi

\iftemplate
\includepdf[pages=-]{../out/\lypdfname.pdf}
\fi

\end{document}